% =====================================================================
% Section 4: Method — EvidenceLink
% =====================================================================
%
% Status: Draft §4 Method v10 (2026-05-19)
%   v10: SYNCHRONIZED WITH evidenceflow PACKAGE MAINLINE (post-1acbfad).
%        - Removed fact-node PPR narrative (Eq.3 deleted).
%        - Removed evidence-need mining as an independent step (Eq.4 / Eq.5
%          deleted; their content is folded into the witness-driven
%          transition expansion since the frozen ETv3 expander internally
%          performs both endpoint-transition and relation-grounded scoring).
%        - Graph nodes are documents; OpenIE facts are cross-document
%          transition WITNESSES, not propagation nodes.
%        - §4.2 narrative replaced: PPR → witness-driven transition
%          expansion that yields a query-local document candidate pool.
%        - §4.3 retains Eq.6 (marginal coverage gain) and Eq.7
%          (preservation-constrained argmax). These match the PCEC
%          native_readout in evidenceflow/native_readout.py.
%
%   Equation count: 7 → 4 (Eq.1 [§3], Eq.2 [§4.2], Eq.6 [§4.3], Eq.7 [§4.3])
%
% =====================================================================

\section{Method: \methodname{}}
\label{sec:method}

\begin{figure}[t]
\centering
\fbox{\parbox{0.92\linewidth}{\centering\vspace{1.5em}
\textit{Overview of \methodname{}.}\\[0.3em]
Indexing: corpus $\to$ source-grounded evidence-link substrate.\\
Retrieval: $\mathcal{S}_q \to G_q \to C_q$ (top-100 documents).\\
Composition: top-$m$ preservation $+$ residual admission $\to R_q^{(K)}$.
\vspace{1.5em}}}
\caption{Overview of \methodname{}. Indexing extracts OpenIE facts and
binds each fact to its source passage as a cross-document transition
witness. At query time, dense entry seeds and question-side anchors
together induce a query-local document graph $G_q$, and a witness-driven
transition expansion produces a ranked candidate pool $C_q$. A
fixed-budget composition step then preserves the top-$m$ entries of
$C_q$ and admits one residual document into the reader prefix
$R_q^{(K)}$ via noisy-OR coverage marginal gain.}
\label{fig:overview}
\end{figure}

We propose \methodname{}, a graph-based retrieval algorithm that
treats query-local graph traversal as the start of evidence selection rather
than the final ranking (Figure~\ref{fig:overview}). \methodname{}
builds a source-grounded evidence-link substrate
(\S\ref{sec:method:substrate}), induces a
query-local document graph and produces a ranked candidate pool
through witness-driven transition expansion
(\S\ref{sec:method:retrieval}), then composes the final reader
prefix under a fixed budget through preservation and residual
admission (\S\ref{sec:method:enhancement}).

\subsection{Source-Grounded Evidence Indexing}
\label{sec:method:substrate}

Our indexing departs from prior OpenIE-graph formulations
\citep{jimenez2024hipporag,gutierrez2025hipporag,wang2025proprag}
in a single move: graph nodes are documents, and each OpenIE fact
is bound to its source passage as a
\emph{cross-document transition witness} rather than as a graph node.
The source-grounded evidence-link substrate is built once for the corpus and reused at query
time without any further OpenIE extraction.

\paragraph{Knowledge Extraction.}
For each document $d \in \mathcal{D}$, a language model extracts
OpenIE facts $f=(h,r,t)$ with argument phrases $h, t$ and relation
phrase $r$
\citep{angeli2015leveraging,gutierrez2025hipporag}. We retain every
fact as a first-class indexing unit and bind it to its source via
$\mathrm{src}(f)=d$, yielding a global fact set $\mathcal{F}$ with
$\mathcal{F}_d = \{f\in\mathcal{F} : \mathrm{src}(f)=d\}$. Retaining
$f$ together with $\mathrm{src}(f)$ keeps the relational triple
intact, so the fact's role is determined by both its argument
structure and its source passage; this provenance binding is what
turns each fact from a free-floating triple into a
document-anchored connection witness.

\paragraph{Graph Construction.}
Documents are connected by two complementary witness families
derived from the same indexing pass. \emph{Endpoint witnesses}
certify a transition between two documents $d, d'$ when there exist
facts $f\in\mathcal{F}_d, f'\in\mathcal{F}_{d'}$ that share at least
one informative argument phrase under
$\{h,t\}\cap\{h',t'\}\neq\emptyset$, encoding cross-document mentions
of the same entity. \emph{Co-source witnesses} group facts within a
document under $\mathrm{src}(f)=\mathrm{src}(f')$, encoding
intra-document co-assertion. Each fact also carries a
\emph{provenance link} to $\mathrm{src}(f)$ used at composition
time. Dense embeddings of fact units and relation phrases are
pre-computed; the source-grounded evidence-link substrate is materialized as an adjacency index
with a hub-degree cap of $30$ and a synonymy threshold of $0.85$ on
dense endpoint similarity (Appendix~\ref{app:hyperparams}).

\subsection{Query-Local Evidence Linking}
\label{sec:method:retrieval}

\paragraph{Query-local document graph.}
Given $q$, \methodname{} first induces a query-local document graph
$G_q$ over the source-grounded evidence-link substrate. We extract entity anchors
$A(q) \subseteq \mathcal{V}_{\mathrm{phr}}$ from $q$ using a
controlled language model that operates on the question only,
never on retrieved documents, and a standard dense retriever
provides an entry signal $\mathrm{Entry}(q) \subseteq \mathcal{D}$.
Let $\mathcal{S}_q$ denote the seed document set, formed by the
union of $\mathrm{Entry}(q)$ and the source documents of facts
whose endpoints overlap $A(q)$. From $\mathcal{S}_q$, $G_q$ is the
$h$-hop induced subgraph along endpoint and co-source witnesses,
\begin{equation}
\label{eq:gq-induced}
G_q = \mathcal{N}_{h}(\mathcal{S}_q),
\end{equation}
which keeps the activated region compact while still admitting
documents reachable through multi-hop entity chains even when their
surface similarity to $q$ is low.

\paragraph{Witness-driven transition expansion.}
$G_q$ alone does not produce a passage ranking. We score documents
in $G_q$ by aggregating witness strengths along the transitions
that connect them to $\mathcal{S}_q$: for each edge $(d, d')$, the
endpoint and co-source witnesses contribute additively, weighted by
the dense relation-phrase similarity $\sigma(\mathbf{r}_f,
\mathbf{q})$ of the certifying fact $f$ and the rank prior of
$\mathcal{S}_q$. The resulting score $s_{\mathrm{link}}(d \mid q)$
is the \emph{evidence-link score} on the document layer: it is local to
$q$ and grounded by the offline witnesses, but evaluated through a
deterministic graph traversal rather than through fact-level
random-walk diffusion. Sorting $G_q$ by $s_{\mathrm{link}}$ yields
the query-local candidate pool $C_q$ of size $|C_q|=100$, on which
the composition step of \S\ref{sec:method:enhancement} operates.
The full traversal procedure, including the hub cap and
relation-phrase weighting, is detailed in
Appendix~\ref{app:graph-details}; we refer to this stage as
\emph{query-local evidence linking} when discussing ablations
(\S\ref{sec:experiments:ablation}).

\subsection{Evidence-Coverage-Aware Retrieval Optimization}
\label{sec:method:enhancement}

The candidate pool $C_q$ is ordered but does not yet match the
reader budget. With a target prefix size $K$ much smaller than
$|C_q|$, two reader-facing problems remain. First, the head of $C_q$
can crowd around a single entity neighborhood, leaving other
question-side requirements uncovered in any short prefix. Second,
strictly trusting the top-$K$ of $C_q$ commits to a passage-level
ranking that does not account for set-level coverage of distinct
requirements. We therefore compose the final reader prefix
$R_q^{(K)}$ through a \emph{preservation-and-admission} readout
that is exact at the default budget $K=5$, $m=4$.

\paragraph{Question-side coverage.}
We extract question-side requirements $B(q)=\{b_1,b_2,\ldots\}$
from $q$ using the same controlled language model used for $A(q)$.
For document $d$ and requirement $b$, $\phi_{d,b}\in[0,1]$ is the
NV-Embed-v2 cosine similarity between $d$ and the requirement
subquery, min--max normalized over $C_q$ and multiplied by the
normalized $d$-to-entity compatibility when $b$ depends on a
resolved bridge entity. Defining
$\mathrm{Cov}_q(R)=\sum_{b\in B(q)}\bigl[1-\prod_{d'\in R}(1-\phi_{d',b})\bigr]$,
the marginal coverage gain of adding $d$ is
\begin{equation}
\label{eq:marginal-gain}
\Delta_q(d \mid R) = \mathrm{Cov}_q(R\cup\{d\}) - \mathrm{Cov}_q(R).
\end{equation}

\paragraph{Preservation and residual admission.}
We form $R_q^{(K)}$ by preserving the top-$m$ entries of $C_q$ as
a fixed prefix $P_q$ and admitting the remaining $K-m$ slots
through greedy coverage-aware selection
\citep{carbonell1998mmr,lin2011submodular} over the residual pool,
at each step adding
\begin{equation}
\label{eq:greedy-rerank}
d^{*}
=
\operatorname*{arg\,max}_{d \in C_q \setminus (P_q \cup R)}
\lambda\,s_{\mathrm{link}}(d \mid q)
+
(1-\lambda)\,\Delta_q(d \mid P_q \cup R),
\end{equation}
with $\lambda\in[0,1]$ tuned once on a development split. At the
default $K=5$, $m=4$ the residual budget is $K-m=1$, so
\eqref{eq:greedy-rerank} reduces to a one-slot exact admission
that either accepts the single residual that maximally improves
coverage or keeps the original $C_q$ top-$K$ unchanged. Larger
residual budgets recover the full greedy procedure.
The reader consumes the resulting prefix of length $K$.
